% =====================================================================
\section{What the model has already failed to do}\label{sec:fail}
The negative results are listed here separately and deliberately. Each of them is a check written down before the computation which did not pass; each narrowed the range of possible explanations more than any of the successes did.

\begin{enumerate}
\item \textbf{The daily range is not a gate but a volume knob.} It was expected that there exists a threshold value of the daily range of the load below which collapses are almost absent and above which they occur. What came out was a smooth growth with no break: over a range from 0 to 40\,\% the fraction of cells that were damaged grew only 1.5-fold. The range changes the risk but switches nothing on.\tD
\item \textbf{Smoothness comes from the spread, not from the network.} The hope was that coupling through network activity would itself create intermediate states between ``healthy'' and ``damaged''. It does not: with a narrow spread of cell resilience the smoothness almost vanishes (fig.~\ref{fig:res}a).\tD
\item \textbf{There is no sleep precursor without a slow variable.} As long as the only source of randomness varied over 2.4 hours, no change of sleep in the day before a collapse arose --- nor could it: the model contained no quantity living on the scale of days. The precursor appeared only after such a quantity was added.\tD
\item \textbf{The episode duration remained wrong.} In the model an episode lasts about six days, whereas observation gives from hours to a few. This is a direct consequence of the law of section~\ref{sec:law}, and it is not corrected by parameters.\tD
\item \textbf{Replacing the ``fraction broken'' by the ``mean damage'' changed nothing.} The argument was strong: the fraction is zero until the first cell breaks, whereas the mean grows at once. In fact, in a narrow window the cells break almost simultaneously, and the two quantities agreed to within half a per cent.\tD
\item \textbf{The serotonin surge turned out to be not a trigger but the same volume knob.} The check was direct: rare one-sided surges raising the load were to cover the upward gap and leave the return to the ordinary noise. Of the rule's four parts not one passed. Collapses following a surge amounted to 20\,\% instead of the expected 80; the episode duration did not change at all --- 5.8 days without surges and 5.2 with them. Surges add collapses (the fraction of affected sets grows from 0.59 to 0.73 as surges become more frequent) but change no property of the episode itself.\tD
\item \textbf{The episode does not remember how deep it went.} What was tested was whether an episode lives a life of its own after the load has returned to normal. It turned out that it almost does not: the episode duration coincides with the time the load spent above the upper boundary of the window, with a rank association of 0.94, whereas the association of the recovery tail with the depth of damage attained is only 0.11. The tail meanwhile occupies a quarter of the episode and lasts about half a day regardless of how many cells broke.\tD
\item \textbf{The first rule turned out to be almost empty.} The check on damage leading the defence passed, but it then emerged that a first-order element satisfies it almost always. A lead of a third of a period, which is observed in experiments, is inaccessible to this core for any parameters.\tD
\end{enumerate}

\begin{note}
Twice during the work a rule passed ``idly'': it was formally satisfied, but on examination it turned out that an extraneous effect was responsible. Once different sets were being compared instead of one and the same; the other time the control rested on a single set, and the statistic simply did not compute. Both cases came to light only because for every rule not only the check but also a way of fooling it had been written down in advance.
\end{note}
